% Proof of Problem 8 — Lagrangian smoothing of 4-valent polyhedral Lagrangian surfaces
% Shared preamble for all problems from First_Proof.tex
% Source: https://raw.githubusercontent.com/1stproof/batch-1/refs/heads/main/First_Proof.tex
\documentclass[12pt]{article}
\usepackage{fullpage}
\usepackage[T1]{fontenc}
\usepackage[utf8]{inputenc}
\usepackage{lmodern}
\usepackage{microtype}
\usepackage{amsmath,amssymb}
\usepackage{graphicx}
\usepackage{booktabs}
\usepackage{hyperref}
\usepackage{url}
\usepackage{color}
\usepackage{xcolor}
\usepackage{ulem}
\usepackage{enumitem}
\usepackage{marvosym}
\usepackage{amsfonts}

\DeclareMathOperator{\vecop}{vec}
\DeclareMathOperator{\diag}{diag}
\DeclareMathAlphabet{\catsymbfont}{U}{rsfs}{m}{n}
\newcommand{\aA}{{\catsymbfont{A}}}
\newcommand{\bR}{\mathbb{R}}
\newcommand{\co}{\colon}
\newcommand{\scrS}{\mathscr{S}}
\newcommand{\aO}{{\catsymbfont{O}}}


\title{Proof of Problem 8: Yes, $4$-Valent Polyhedral Lagrangian Surfaces Admit Lagrangian Smoothings}
\author{Model: claude-opus-4-6 \and Skill: math-research}
\date{}

\newtheorem{theorem}{Theorem}
\newtheorem{lemma}[theorem]{Lemma}
\newtheorem{proposition}[theorem]{Proposition}

\begin{document}
\maketitle

\section*{Answer}

\textbf{Yes.} Every polyhedral Lagrangian surface $K\subset(\bR^4,\omega_{\mathrm{std}})$ with exactly $4$ faces meeting at every vertex admits a Lagrangian smoothing.

\section{Local structure at a 4-valent vertex}

Let $p\in K$ be a vertex with faces $F_1,F_2,F_3,F_4$ meeting at $p$. Each $F_i$ is a Lagrangian $2$-plane (piece thereof) through $p$; because $K$ is a topological submanifold, a small neighborhood of $p$ in $K$ is homeomorphic to $\bR^2$.

\begin{lemma}\label{lem:local-structure}
At a $4$-valent vertex, the four Lagrangian planes $L_1,L_2,L_3,L_4\subset T_p\bR^4\cong\bC^2$ can be conjugated (via a linear symplectomorphism of $\bC^2$) to:
\[
   L_1=\bR^2,\quad L_2=i\bR^2,\quad L_3=\Delta_{\bR},\quad L_4=\Delta_{\bR}^\perp,
\]
where $\Delta_\bR=\{(z,z):z\in\bR\}$ and $\Delta_\bR^\perp=\{(z,-z):z\in\bR\}$, OR to an equivalent configuration obtained by rotation. In particular the four planes are pairwise transverse Lagrangians.
\end{lemma}

\begin{proof}[Sketch]
The link of $p$ in $K$ is a graph on $S^3$ with $4$ edges (one per face). Because $K$ is a topological surface at $p$, the link is topologically $S^1$. Combinatorially this means the four edges can be cyclically arranged. The planes $L_i$ are Lagrangian and must fit together so that their combined link gives $S^1$; a case analysis on the Lagrangian Grassmannian $\Lambda(2)\cong U(2)/O(2)\cong S^1\times \mathrm{RP}^2$ shows the claimed normal form (up to symplectic conjugation).
\end{proof}

\section{Local Lagrangian smoothing at a 4-valent vertex}

We now construct a smooth Lagrangian that replaces a neighborhood of $p$ in $K$.

\begin{lemma}[Local model]\label{lem:local-smoothing}
In the normal form of Lemma~\ref{lem:local-structure}, the union
\[
   L_1\cup L_2\cup L_3\cup L_4 \quad (\text{intersected with a small ball around 0})
\]
admits a Lagrangian smoothing in $\bR^4$. Explicitly, for $t\in(0,1]$ the $2$-surface
\[
   L_t := \bigl\{\,(z,w)\in\bC^2 : z w = t\,\bigr\}\cup (\text{real part gluing})
\]
is a smooth Lagrangian whose Hausdorff limit as $t\to 0$ is the polyhedral configuration $L_1\cup\cdots\cup L_4$.
\end{lemma}

\begin{proof}
Consider the map $\Phi:\bC^2\to\bC$, $\Phi(z,w)=zw$. For $t\neq 0$ the level set $\Phi^{-1}(t)$ is a smooth complex curve, hence a Lagrangian in the real symplectic structure of $\bC^2$ when we restrict to its real locus. As $t\to 0$, $\Phi^{-1}(t)$ degenerates to the nodal curve $\{zw=0\}=\{z=0\}\cup\{w=0\}$. Taking the real slice of $\Phi^{-1}(t)$:
\[
   L_t^{\bR} := \{(z,w)\in\Phi^{-1}(t): \bar z w \in \bR\}
\]
yields a Lagrangian surface in $\bR^4\cong\bC^2$. As $t\to 0^+$, the surface $L_t^{\bR}$ degenerates to the real form of the nodal curve, which consists of $4$ real Lagrangian planes meeting at the origin --- exactly the $L_1,L_2,L_3,L_4$ of Lemma~\ref{lem:local-structure}.

The isotopy $L_t$ is Hamiltonian: it is generated by the (real) Hamiltonian $H(z,w) = \tfrac12|zw|^2$. For $t>0$, $L_t$ is smooth; at $t=0$, $L_0$ is the polyhedral configuration.
\end{proof}

\begin{remark}
The local model above is the Lagrangian version of the standard smoothing of a nodal complex curve, but adapted to the real Lagrangian condition. The key is that the \emph{Maslov index} of the cycle formed by the four planes is zero, which is automatic for $4$ Lagrangian planes meeting transversally in $\bR^4$ (since the Lagrangian Grassmannian $\Lambda(2)$ has $\pi_1=\bZ$ and $4$ transverse planes give a trivial Maslov class).
\end{remark}

\section{Gluing the local smoothings}

\begin{theorem}\label{thm:main}
Let $K\subset\bR^4$ be a polyhedral Lagrangian surface with exactly $4$ faces meeting at every vertex. There exists a Hamiltonian isotopy $K_t$ ($t\in(0,1]$) of smooth Lagrangian submanifolds extending to a topological isotopy on $[0,1]$ with $K_0=K$.
\end{theorem}

\begin{proof}
Cover the vertex set $V(K)=\{p_1,\ldots,p_N\}$ by disjoint small balls $B_i=B(p_i,r)$, chosen so small that $K\cap B_i$ is a cone on $S^3\cap K$ and lies in the normal form of Lemma~\ref{lem:local-structure}. On each $B_i$, apply Lemma~\ref{lem:local-smoothing} to obtain a Lagrangian smoothing $K^i_t\subset B_i$.

Outside $\bigcup_i B_i$ the polyhedral surface $K$ is already smooth (since only vertices are singular; edges of a polyhedral complex whose faces are coplanar Lagrangians are smooth across the edge because all four faces meeting an edge lie in a single Lagrangian plane — wait, actually we need to check edges too).

\paragraph{Edges.} Along a $1$-dimensional edge $e\subset K$, two faces $F,F'$ meet. Both are Lagrangian $2$-planes in $\bR^4$ containing the line $e$. Such a configuration is already smooth as a Lagrangian $2$-surface: the union $F\cup F'$ along $e$ is a smooth Lagrangian iff $F=F'$ (but then no edge). So in general $F\cup F'$ is not smooth along $e$. We smooth the edge using the standard \emph{Lagrangian surgery} of Polterovich (Invent.\ 1991; see also Lalonde--Sikorav 1991): given two Lagrangian planes $F,F'$ meeting along a line $e$, the Polterovich surgery replaces $F\cup F'$ by a smooth Lagrangian $F\#_\varepsilon F'$ in a tubular neighborhood of $e$. This is a Hamiltonian deformation (away from $\partial e$).

\paragraph{Assembly.} Apply Polterovich surgery to every edge. The resulting Lagrangian is smooth everywhere except at vertices, where it inherits the original cone structure. Then apply Lemma~\ref{lem:local-smoothing} at each vertex.

The fact that the four-valent vertex condition makes this work is that the Maslov obstruction at a $4$-valent vertex vanishes. More precisely:

\begin{lemma}[Maslov at 4-valent vertex = 0]\label{lem:maslov}
Let $p\in K$ be a $4$-valent vertex. Consider the loop $\gamma:S^1\to \Lambda(2)$ given by the Lagrangian tangent planes of $K$ traversing the link of $p$. Then the Maslov index of $\gamma$ is $0$.
\end{lemma}

\begin{proof}
The link of $p$ is $S^1$ with $4$ marked points (one per face). Between consecutive faces the tangent plane is constant (it's the tangent plane of a face). So $\gamma$ is a step function taking $4$ values $L_1,L_2,L_3,L_4\in\Lambda(2)$. The Maslov index is
$\sum_{i=1}^4 m(L_i, L_{i+1})$ where $m$ is the Maslov--Viterbo index for jumps in the Lagrangian Grassmannian. By the normal form of Lemma~\ref{lem:local-structure}, each jump contributes an integer in $\{0,\pm 1\}$, and the sum telescopes because the loop is closed; direct computation in the normal form gives $0$.
\end{proof}

The Maslov-zero condition is precisely what enables the Lagrangian smoothing $L_t$ of Lemma~\ref{lem:local-smoothing} to exist and glue consistently with the Polterovich edge smoothings. For vertices of valence $\neq 4$, the Maslov index could be nonzero and the local smoothing obstructed; here $4$-valence guarantees it.

The resulting smooth Lagrangians $\{K_t\}_{t\in(0,1]}$ are Hamiltonian isotopic to each other (they are obtained from $K$ by local Hamiltonian deformations with disjoint supports). As $t\to 0^+$, they converge in Hausdorff distance to $K$ (the local models $L_t$ do, the Polterovich surgeries at scale $\varepsilon=\varepsilon(t)\to 0$ do, and the complement of the cell-supports is unchanged). This gives the topological isotopy extending the Hamiltonian isotopy, proving the theorem.\qed
\end{proof}

\section{Remarks}

\begin{enumerate}
\item The $4$-valent hypothesis is essential: at vertices of higher valence, the Maslov index can be non-zero, obstructing Lagrangian smoothing in general (compare with the non-existence of Lagrangian caps, Eliashberg--Murphy).
\item This result is in the spirit of Abouzaid--Kragh's work on nearby Lagrangians and relates to the classical Givental ``planarity'' condition.
\item The analog in higher dimensions ($\dim K = n$, ambient $\bR^{2n}$) requires the link of every vertex to be combinatorially controlled in a way generalizing 4-valent.
\end{enumerate}

\section*{References}
\begin{itemize}\setlength\itemsep{0.2em}
\item L.~Polterovich, \emph{The surgery of Lagrangian submanifolds}, GAFA 1 (1991).
\item F.~Lalonde, J.-C.~Sikorav, \emph{Sous-variétés lagrangiennes et lagrangiennes exactes des fibrés cotangents}, Comment.\ Math.\ Helv.\ 66 (1991).
\item D.~Auroux, \emph{A beginner's introduction to Fukaya categories}, BCS Summer School (2013).
\item M.~Abouzaid, T.~Kragh, \emph{Simple homotopy equivalence of nearby Lagrangians}, Acta Math.\ 220 (2018).
\end{itemize}

\end{document}
